\section{Concurrency Design \& Implementation}
\label{sec:concurrency}

Phase 2 introduces a high-concurrency operating condition: at the start of a
semester, many users may request popular spaces for overlapping periods within
a short interval. Eligible requests are confirmed immediately through instant
booking, whereas other requests follow the staff-approval workflow. Since both
pathways can confirm the same space--time interval concurrently, the database
must enforce NR~6: \emph{two approved bookings must never use the same space
during overlapping periods, regardless of how they were created}.

Task~08 identified five concurrency conflicts (K1--K5) that can violate this
invariant or expose stale availability. The solution was specified in
\texttt{11-concurrency-design}, implemented in
\texttt{12-concurrency-implementation.sql}, and verified through deterministic
two-session tests in \texttt{13-concurrency-tests}. It serializes the four
write paths that can change the confirmed state of a space--time interval:
instant submission, staff approval, maintenance escalation or downgrade, and
maintenance-ticket creation. Each path uses the same logical resource for a
space and rechecks the relevant invariants immediately before writing.

\subsection{Identified Concurrency Conflicts}

The primary conflict is the K1 \emph{read--check--write race}. Two concurrent
requests for the same space and period can both observe the interval as
available and then confirm their bookings before either transaction sees the
other's result. Table~\ref{tab:race-sequence} illustrates this race for two
instant-booking sessions targeting space $S_1$ from 10:00 to 12:00.

\begin{table}[htbp]
    \centering
    \small
    \renewcommand{\arraystretch}{1.15}
    \begin{tabularx}{\textwidth}{c
        >{\RaggedRight\arraybackslash}X
        >{\RaggedRight\arraybackslash}X}
        \toprule
        \textbf{Step} & \textbf{Session A: instant submission} &
        \textbf{Session B: instant submission} \\
        \midrule
        1 & Check $S_1$: no confirmed overlap or blocking maintenance. & \\
        2 & & Check $S_1$: no confirmed overlap or blocking maintenance. \\
        3 & Insert a pending booking and its automatic approval
        (\texttt{approver\_id = -1}). & \\
        4 & & Insert a pending booking and its automatic approval; BR1 and BR4
        checks see no committed competitor. \\
        5 & Commit. & \\
        6 & & Commit. \\
        \bottomrule
    \end{tabularx}
    \caption{Concurrent instant-booking race for the same space and period}
    \label{tab:race-sequence}
\end{table}
\FloatBarrier

At step~4, session A's write remains uncommitted and is therefore not visible
to session B under the default \texttt{READ COMMITTED} isolation level. Both
transactions pass the availability rule (BR1) and maintenance rule (BR4), then
commit overlapping approved bookings on $S_1$, violating BR1 and NR6.

The interval between the availability check and the confirmation write is the
\emph{race window}. Because no lock spans this interval, neither existing
backstop fully closes it. The trigger
\texttt{trg\_bookings\_prevent\_overlap} cannot validate against a competing
row that is not yet visible, while the filtered unique index
\texttt{uq\_bookings\_active\_overlap} prevents only identical start times.
It does not reject intersecting intervals with different start times, such as
10:00--12:00 and 11:00--13:00.

The same check--write race occurs across other booking and maintenance paths.
K2 covers instant submission racing with staff approval. K3 occurs when a
maintenance record is escalated to \texttt{out-of-service} during an in-flight
confirmation, while K5 covers insertion of a new maintenance ticket whose
default impact is \texttt{out-of-service}. K4 is a read-side conflict in which
room-finder availability becomes stale before the user acts.
Table~\ref{tab:conflicts} summarizes the complete conflict inventory.

\begin{table}[htbp]
    \centering
    \small
    \renewcommand{\arraystretch}{1.15}
    \begin{tabularx}{\textwidth}{c
        >{\RaggedRight\arraybackslash}p{0.37\textwidth}
        >{\RaggedRight\arraybackslash}p{0.16\textwidth}
        >{\RaggedRight\arraybackslash}X}
        \toprule
        \textbf{ID} & \textbf{Concurrent scenario} &
        \textbf{Affected rules} & \textbf{Path} \\
        \midrule
        K1 & Two requests read the same interval as free and both confirm. &
        BR1, NR6 & Booking insertion \\
        K2 & Instant submission races with staff approval. &
        BR1, NR6 & Cross-pathway confirmation \\
        K3 & Maintenance is escalated while a booking confirmation is in flight. &
        BR4, NR4 & Maintenance update \\
        K4 & Room-finder availability becomes stale before confirmation. &
        BR1, NR4, NR6 & Read-only availability \\
        K5 & An out-of-service ticket is inserted during confirmation. &
        BR4, NR6 & Maintenance insertion \\
        \bottomrule
    \end{tabularx}
    \caption{Concurrency conflicts identified in the Task 08 analysis}
    \label{tab:conflicts}
\end{table}
\FloatBarrier

K1, K2, K3, and K5 share the same root cause: concurrent transactions perform
a check and a conflicting write on the same space--time interval without a
common serialization point. K4 cannot by itself corrupt the database, but its
stale result must not be trusted during confirmation. The implemented control
therefore serializes the check--recheck--write sequence per space and validates
availability again inside the protected transaction. Consequently, two
conflicting confirmations cannot both commit.

\subsection{Cause and Chosen Control}

\subsubsection{Cause}

The race occurs under the default READ COMMITTED isolation level
because a transaction cannot observe another transaction's uncommitted rows.
Two writers can therefore validate the same availability state before either
commits. The existing triggers execute within their writer's transaction and
cannot validate against an invisible competitor, while the filtered unique
index rejects only identical start times. Since no common serialization point
spans the check--write interval, all confirmation-related write paths require
one.

\subsubsection{Application-Lock Protocol}

The implementation uses a transaction-owned SQL Server application lock. After
starting a transaction, each write procedure calls
\textbf{sys.sp\_getapplock} with the following configuration:

\begin{itemize}
    \item Resource: \textbf{space\_booking:<space\_id>};
    \item Mode: \textbf{Exclusive};
    \item Owner: \textbf{Transaction};
    \item Timeout: 5 seconds.
\end{itemize}

The lock is released automatically when the transaction commits or rolls back.
The procedures \textbf{usp\_booking\_instant\_submit},
\textbf{usp\_booking\_approve},
\textbf{usp\_maintenance\_set\_impact\_level}, and
\textbf{usp\_maintenance\_report} request the same resource for a given space.
This serializes instant approvals, staff approvals, and booking--maintenance
interactions through one protocol. Maintenance reporting acquires the lock only
when the new ticket begins as \textbf{out-of-service}; an advisory ticket does
not block the interval and instead creates an acknowledgement requirement.

\subsubsection{Invariant Revalidation and Failure Handling}

The lock protects the critical section, but correctness depends on revalidation
inside it. Immediately before writing, each procedure rereads its target data
and checks confirmed-booking overlap (BR1), out-of-service maintenance overlap
(BR4), space status (BR2), capacity (BR3), and acknowledgement completeness
(NR2). A failed check rolls back the transaction and returns a deterministic
result code. Existing triggers and the filtered unique index remain as
defense-in-depth controls.

Lock-acquisition failures are returned to the caller rather than silently
retried. A timeout returns \texttt{51005}, a cancelled request returns
\texttt{51006}, and both are retryable. A deadlock victim returns
\texttt{51007}; in this case, the caller must restart the entire transaction.

\subsubsection{Design Rationale and Trade-offs}

The application lock is valid under READ COMMITTED and requires no
schema change. In contrast, SERIALIZABLE key-range locking or UPDLOCK / HOLDLOCK range reads are reliable only when every
writer locks the same predicate range. The four write paths use different
interval operations, making this approach plan-sensitive and vulnerable to a
range-boundary error that could reopen K3 or K5. Optimistic concurrency would
require version columns excluded from the approved design, while triggers alone
cannot serialize concurrent writers under READ COMMITTED.

The lock resource is keyed by space rather than by space and day. A daily key
would require a multi-day booking to acquire several locks in a fixed order,
creating additional lock-ordering and cross-midnight risks. The per-space key
avoids these cases but serializes all writers for a popular space. This cost is
acceptable for the expected workload: Task~14 generates more than 100{,}000
bookings across three academic years, corresponding to tens of writes per day
for the busiest space, and each critical section completes well within the
five-second timeout. If production monitoring later identifies material
contention, a per-day resource key remains the documented tuning option.

\subsection{Concurrent Test Results}
Document the test setup, session ordering, expected outcome, and actual result.
Include concise logs or screenshots demonstrating both the unsafe behavior
before the fix and the prevented conflict after the fix.

\begin{table}[htbp]
    \centering
    \begin{tabularx}{\textwidth}{p{0.19\textwidth}X X X}
        \toprule
        \textbf{Test} & \textbf{Concurrent actions} & \textbf{Before fix} & \textbf{After fix} \\
        \midrule
        % TODO: Populate from 13-concurrency-tests.
        To be completed & To be completed & To be completed & To be completed \\
        \bottomrule
    \end{tabularx}
    \caption{Concurrency test results}
    \label{tab:concurrency-tests}
\end{table}